\documentclass[10pt,a4paper]{article}
\usepackage{amsmath,amssymb,amsthm,graphicx}
\usepackage{amsmath,amssymb,amsthm,hyperref}
\newtheorem{nttheorem}{Theorem}
\newtheorem{ntdefinition}[nttheorem]{Definition}
\newtheorem{ntlemma}[nttheorem]{Lemma}
\newtheorem{ntcorollary}[nttheorem]{Corollary}
\title{Hardness of approximating deterministic communication complexity of partial Boolean functions}
\author{Rahul Ilango, Bruno Loff, Igor C. Oliveira}
\date{doi:10.4230/LIPIcs.CCC.2020.22; CC BY3.0}
\begin{document}
\maketitle
\noindent\textit{Editorial note.} Rahul Ilango,Bruno Loff and IgorC.Oliveira. CCC2020 article22,DOI10.4230/LIPIcs.CCC.2020.22. CC BY3.0.\par\medskip
\noindent\textit{Editorial note.} Faithful editable transcription of the official CCC2020 CC BY3.0 publication, article22. Independent partial-Boolean communication-minimization theme: complete Section5.1 source definitions19--31 and cited coloring-hardness20/21 premises, complete Section5.2 graph-to-partial-matrix construction and Theorem32 author proof, complete Section5.3 original Corollaries33/34 with author deductions and Theorem35 full approximate-characterization impossibility proof. All82 official bibliography entries retained using the shared frozen formal-source bibliography. Original partial\{0,1,*\} matrix, unrestricted deterministic two-party protocol/leaf/depth/cover/DAG measures, approximation ratios/error and eventual monotonicity/gap conditions retained. No fixed-round or total-function theorem substituted. Original printed references/wording retained without mathematical repair, including Section5.3 reference to balanced protocol as proofTheorem20 and source DAG counting statement. Standard external graph-coloring hardness retained at author citation boundary, without recursive proof. All direct parameter consequences/communication measures grouped once. Finite comparison to NT065: that item uses CCC2021.31 deterministic totalBoolean fixed-round protocol-size hardness with a Boolean construction and round elimination; this item uses CCC2020.22 partialBoolean matrices, arbitrary depth and coloring-cover characterization. Distinct model, main result and mechanism. Native not available in bounded author/source lookup; all selected TeX transcribed from formally licensed PDF. No necessary figures. Root handles compilation/database/publication.\par\medskip
\noindent\textit{Allowed standard prerequisites (editorial).}\begin{itemize}
\item Gap hardness of graph coloring. Exact original Theorem20[82] and Theorem21[79] statements retained; no recursion into prior hardness proofs.
\end{itemize}\medskip
\setcounter{nttheorem}{34}
\begin{nttheorem}[Author Theorem35]
Let $\lambda,\eta:\mathbb R_{\geq0}\to\mathbb R_{\geq0}$ be functions such that $\eta(1+\log x)<\lambda(0.99x)$ for all sufficiently large $x$. Moreover, assume that these functions are non-decreasing for every large enough $x$. Then, if $\operatorname P\ne\operatorname{NP}$, there is no polynomial-time computable function $r$, which accepts as input the communication matrix $M_f\in\{0,1,*\}^{n\times n}$ of a partial Boolean function $f:[n]\times[n]\to\{0,1,*\}$, and for every large enough $n$ outputs a value $r(M_f)$ such that
\[\lambda(D(f))\leq r(M_f)\leq\eta(D(f)).\]
\end{nttheorem}

\section*{5 On the NP-hardness of communication minimization problems}
\subsection*{5.1 Background}
\paragraph{Hardness of graph coloring}
Our NP-hardness reduction is from the chromatic number problem:
\setcounter{nttheorem}{18}
\begin{ntdefinition}[Author Definition19: Chromatic number]
A coloring of an undirected graph $G$, is a partition of the vertices such that no edge has both endpoints in the same part. The chromatic number of a graph $G$, denoted $\chi(G)$, is the smallest number of parts of a coloring of $G$.
\end{ntdefinition}
The NP-hardness of approximating the chromatic number has been established by a series of results~\cite{ref57,ref31,ref24}, culminating in a paper by Zuckerman~\cite{ref82}, where the following was proven:
\begin{nttheorem}[Author Theorem20, {\cite{ref82}}]
For every constant $\epsilon>0$ it is NP-hard to approximate $\chi(G)$ for a given $n$-vertex graph $G$, with an approximation ratio better than $n^{1-\epsilon}$. More precisely, for every $L\in\operatorname{NP}$ and $\epsilon>0$, there exists a polynomial-time algorithm that, on input $x$, outputs a parameter $k$ and an $n$-vertex graph $G$ such that if $x\in L$ then $\chi(G)\leq k$, and if $x\notin L$ then $\chi(G)>n^{1-\epsilon}\cdot k$.
\end{nttheorem}
In the reductions above, the parameter $k$ is $\Theta(n^\epsilon)$. More recent results on the hardness of approximating chromatic number allow for a different gap, where $k$ is constant as $n$ grows, and we wish to distinguish graphs which are $k$-colorable from graphs which are not $g(k)$-colorable, for a fast-growing function $g:\mathbb N\to\mathbb N$. An original such result with $g(k)=k^{\Omega(\log k)}$ was shown by Khot~\cite{ref47}, which was later improved to $g(k)=2^{\Omega(k^{1/3})}$ by Huang~\cite{ref37}. The latest result, by Wrochna and \v Zivn\'y~\cite{ref79}, achieves $g(k)=\binom k{\lfloor k/2\rfloor}$:
\begin{nttheorem}[Author Theorem21, {\cite{ref79}}]
Let $k\geq4$ be a natural number. For every $L\in\operatorname{NP}$, there exists a polynomial-time algorithm that, on input $x$, outputs an graph $G$ such that if $x\in L$ then $\chi(G)\leq k$, and if $x\notin L$ then $\chi(G)>\binom k{\lfloor k/2\rfloor}$.
\end{nttheorem}
\paragraph{Communication complexity of relations and partial functions}
We will need the following definitions.
\begin{ntdefinition}[Author Definition22]
We will call a two-player communication relation, or simply a relation, to any subset $F\subseteq\mathcal X\times\mathcal Y\times\mathcal Z$, where $\mathcal X,\mathcal Y,\mathcal Z$ are finite sets, such that, for every $(x,y)\in\mathcal X\times\mathcal Y$, there exists at least one $z\in\mathcal Z$ with $(x,y,z)\in F$.
\end{ntdefinition}
Given a relation $F\subseteq\mathcal X\times\mathcal Y\times\mathcal Z$, and a pair $(x,y)\in\mathcal X\times\mathcal Y$, we let
\[F(x,y)=\{z\in\mathcal Z\mid(x,y,z)\in F\}\ne\varnothing.\]
The matrix with $\mathcal X$-indexed rows and $\mathcal Y$-indexed columns, whose $(x,y)$ entry is $F(x,y)$, is called the communication matrix of $F$.
Given a relation $F\subseteq\mathcal X\times\mathcal Y\times\mathcal Z$, a rectangle $R=A\times B\subseteq\mathcal X\times\mathcal Y$, and an element $z\in\mathcal Z$, we say $R$ is $z$-monochromatic for $F$, if $(x,y,z)\in F$ for all $(x,y)\in R$. We say $R$ is monochromatic for $F$ if there exists a $z\in\mathcal Z$ with $R$ being $z$-monochromatic for $F$.
A partial two-player function is a relation $f\subseteq\mathcal X\times\mathcal Y\times\mathcal Z$ such that, for every $(x,y)\in\mathcal X\times\mathcal Y$, either $|f(x,y)|=1$, in which case we say $f$ is defined at $(x,y)$, or $f(x,y)=\mathcal Z$, in which case we say $f$ is undefined at $(x,y)$. If $f$ is defined at $(x,y)$, we write $f(x,y)=z$ in place of $f(x,y)=\{z\}$, and if $f$ is undefined at $(x,y)$, we write $f(x,y)=*$, instead of $f(x,y)=\mathcal Z$.
\begin{ntdefinition}[Author Definition23: $P(F)$]
Given a relation $F\subseteq\mathcal X\times\mathcal Y\times\mathcal Z$, a partition of $F$ is a family $\Pi$ of monochromatic rectangles for $F$, such that every $(x,y)\in\mathcal X\times\mathcal Y$ appears in exactly one rectangle of $\Pi$. The partition number of a relation $F\subseteq\mathcal X\times\mathcal Y\times\mathcal Z$, denoted $P(F)$ is the smallest possible size of a partition of $F$.
\end{ntdefinition}
\begin{ntdefinition}[Author Definition24: $C(F)$]
Given a relation $F\subseteq\mathcal X\times\mathcal Y\times\mathcal Z$, a cover of $F$ is a family $\Gamma$ of monochromatic rectangles for $F$, such that every $(x,y)\in\mathcal X\times\mathcal Y$ appears in at least one rectangle of $\Gamma$. The cover number of a relation $F\subseteq\mathcal X\times\mathcal Y\times\mathcal Z$, denoted $C(F)$ is the smallest possible size of a cover of $F$.
\end{ntdefinition}
In the context of Boolean-valued functions, the usual notion of cover number (see~\cite{ref42} Section4.2,p97), and related notion of non-deterministic communication complexity, only require that the 1s of the communication matrix are covered.
\begin{ntdefinition}[Author Definition25: $C_1(F)$]
Given a partial two-player function $f:\mathcal X\times\mathcal Y\to\{0,1\}$, a 1-cover of $F$ is a family $\Gamma$ of 1-monochromatic rectangles for $F$, such that every $(x,y)\in\mathcal X\times\mathcal Y$ having $f(x,y)=1$ appears in at least one rectangle of $\Gamma$. The1-cover number of $f$, denoted $C_1(F)$ is the smallest possible size of a 1-cover of $F$. We then define the non-deterministic communication complexity of $f$ to be $N(f)=\lceil\log C_1(f)\rceil$.
\end{ntdefinition}
\begin{ntdefinition}[Author Definition26]
A protocol $\pi$ over $\mathcal X\times\mathcal Y$ is a rooted tree:
\begin{itemize}
\item Each node $v$ is associated with a rectangle $\pi^{-1}(v)=A\times B\subseteq\mathcal X\times\mathcal Y$.
\item Each non-leaf node $v$, with $\pi^{-1}(v)=A\times B$, is labeled by either(a) a partition $A=A_0\mathbin{\dot\cup}A_1$ of $A$, in which case we say it is Alice's node or(b) a partition $B=B_0\mathbin{\dot\cup}B_1$ of $B$, in which case we say it is Bob's node.
\item The rectangle associated with the root is $\mathcal X\times\mathcal Y$.
\item If a non-leaf node $v$ of Alice has $\pi^{-1}(v)=A\times B$ and is labeled by a partition $A=A_0\mathbin{\dot\cup}A_1$ of $A$, then for each $c\in\{0,1\}$ there will be a unique child $v_c$ of $v$, with $\pi^{-1}(v_c)=A_c\times B$; similarly for Bob's nodes.
\end{itemize}
\end{ntdefinition}
\begin{ntdefinition}[Author Definition27: $L(F),D(F)$]
Let $F\subseteq\mathcal X\times\mathcal Y\times\mathcal Z$ be a relation, and $\pi$ a protocol over $\mathcal X\times\mathcal Y$. We say that $\pi$ is a protocol for $F$ if, for any leaf $\ell$ of $\pi$, the rectangle $\pi^{-1}(\ell)$ is monochromatic for $F$. We then let $L(F)$ be the smallest possible number of leaves in a protocol for $F$, and we let $D(F)$ be the smallest possible depth of a protocol for $F$.
\end{ntdefinition}
Note that the rectangles associated with the leaves of a protocol for $F$ form a partition of $F$, any partition of $F$ is a cover of $F$, and any cover of $F$ contains a 1-cover of $F$, and hence we get the following:
\begin{ntlemma}[Author Lemma28]
For any relation $F\subseteq\mathcal X\times\mathcal Y\times\mathcal Z$, $L(F)\geq P(F)\geq C(F)$, and for a partial two-player function $f:\mathcal X\times\mathcal Y\to\{0,1\}$, $C(f)\geq C_1(f)$.
\end{ntlemma}
\begin{ntdefinition}[Author Definition29]
A DAG-like protocol $\gamma$ over $\mathcal X\times\mathcal Y$ is a directed acyclic graph:
\begin{itemize}
\item $\gamma$ has a single source node, called the root, and one or more sink nodes, called leaves.
\item Each node $v$ is associated with a rectangle $\gamma^{-1}(v)=A\times B\subseteq\mathcal X\times\mathcal Y$.
\item Each non-leaf node $v$, with $\gamma^{-1}(v)=A\times B$, is labeled by either(a) a partition $A=A_0\mathbin{\dot\cup}A_1$ of $A$, in which case we say it is Alice's node or(b) a partition $B=B_0\mathbin{\dot\cup}B_1$ of $B$, in which case we say it is Bob's node.
\item The rectangle associated with the root is $\mathcal X\times\mathcal Y$.
\item If a non-leaf node $v$ of Alice has $\gamma^{-1}(v)=A\times B$ and is labeled by a partition $A=A_0\mathbin{\dot\cup}A_1$ of $A$, then for each $c\in\{0,1\}$ there will be an edge from $v$ to a node $v_c$ of $\gamma$, which must be such that $A_c\times B\subseteq\gamma^{-1}(v_c)$; similarly for Bob's nodes.
\end{itemize}
\end{ntdefinition}
\begin{ntdefinition}[Author Definition30: $S(F)$]
Let $F\subseteq\mathcal X\times\mathcal Y\times\mathcal Z$ be a relation, and $\gamma$ a DAG-like protocol over $\mathcal X\times\mathcal Y$. We say that $\gamma$ is a DAG-like protocol for $F$ if, for any leaf $\ell$ of $\gamma$, the rectangle $\gamma^{-1}(\ell)$ is monochromatic for $F$. We then let $S(F)$ be the smallest possible number of nodes in a DAG-like protocol for $F$.
\end{ntdefinition}
Note that:(1) a protocol with $k$ leaves is a DAG-like protocol with $2k-1$ nodes,(2) the rectangles associated with the leaves of a DAG-like protocol form a cover of $F$, and(3) any DAG with 1 source node,$k$ sink nodes, and maximum out-degree 2 has at least $2k-1$ nodes in total. Hence:
\begin{ntlemma}[Author Lemma31]
$S(F)\leq2L(F)-1$ and $S(F)\geq2C(F)-1$.
\end{ntlemma}

\subsection*{5.2 A reduction from Graph Coloring to Partial-MCCP}
Let us be given a graph $G=([n],E)$ with $E\subseteq\binom{[n]}2$, and consider the partial two-player Boolean function $f_G:[n]\times[n]\to\{0,1,*\}$, given by
\[f_G(i,j)=\begin{cases}1&\text{if }i=j,\\0&\text{if }\{i,j\}\in E,\\ *&\text{if }\{i,j\}\notin E.\end{cases}\]
By $f_G(i,j)=*$ we mean that $f_G(i,j)$ is undefined, which is to say Alice and Bob may output any value when given $(i,j)$ as input.
\setcounter{nttheorem}{31}
\begin{nttheorem}[Author Theorem32]
We have $L(f_G)\leq2\chi(G)$ and $C_1(f_G)=\chi(G)$.
\end{nttheorem}
\begin{proof}
To show that $L(f_G)\leq2\chi(G)$, consider the simple protocol where Alice and Bob have agreed on a coloring of $G$, so Alice begins by sending the color of her vertex, and Bob replies whether the color of his vertex is the same. If the two colors match, they output 1, and otherwise they output 0.
This is a protocol for $f$, since a valid coloring of $G$ will only color $i$ and $j$ by the same color if $i=j$, or $\{i,j\}\notin E$, in which case 1 is a valid output. If $i$ and $j$ are colored differently, then $i\ne j$, and so 0 is a valid output. This protocol has $2\chi(G)$ leaves.
To show that $C_1(f_G)\leq\chi(G)$, we cover the diagonal by 1-monochromatic rectangles; each such rectangle corresponds to a color $c$, and is the smallest rectangle containing all diagonal entries $(i,i)$ for vertices $i$ of $G$ that are colored $c$.
Now suppose we have any positive cover $\Gamma$ for $f$. The diagonal entries must be covered by 1-monochromatic rectangles of $\Gamma$, so consider the coloring of $G$ which colors vertex $i\in[n]$ by the 1-monochromatic rectangle which covers $(i,i)$. By our choice of $f$, any two diagonal entries belonging to the same 1-monochromatic rectangle are not connected by an edge of $G$. And so this gives us a valid coloring of $G$, which implies that $\chi(G)\leq|\Gamma|$.
\end{proof}
\subsection*{5.3 Consequences of the reduction}
We may now take Theorems32 and20, together with Lemmas28 and31, to obtain:
\begin{ntcorollary}[Author Corollary33]
For any constant $\epsilon>0$, it is NP-hard to approximate $L(f),P(f),C(f),C_1(f)$ and $S(f)$, for a given partial function $f:[n]\times[n]\to\{0,1,*\}$, with an approximation ratio better than $n^{1-\epsilon}$.
\end{ntcorollary}
Observe that this hardness of approximation result with a ratio of $n^{1-\epsilon}$ is very close to optimal, since all aforementioned measures for a partial two-player function $f:[n]\times[n]\to\{0,1,*\}$ are upper-bounded by $2n$. Improving upon this might still be possible, and we leave it as an open problem.
Now notice that the protocol given in the proof of Theorem20 is balanced. For this reason, its depth is logarithmic in the number of leaves. Thus, since computing $L(f)$ is hard up to an approximation ratio of $n^{1-\epsilon}$, it follows that computing the communication complexity $D(f)$ for a given partial function $f:[n]\times[n]\to\{0,1,*\}$ is hard, even if we allow for an additive error term of $(1-\epsilon)\log n$. Since $D(f)\leq\log n$, this implies that we cannot approximate $D(f)$ with an approximation ratio better than $\approx1/\epsilon$, which we can take as an arbitrarily large constant. The same reasoning applies to non-deterministic communication complexity.
\begin{ntcorollary}[Author Corollary34]
For any constant $\epsilon>0$, it is NP-hard to approximate $D(f)$ or $N(f)$, for a given partial function $f:[n]\times[n]\to\{0,1,*\}$, with an error term smaller than $(1-\epsilon)\log n$. For any constant $c>1$, it is NP-hard to approximate $D(f)$ or $N(f)$, for a given partial function $f:[n]\times[n]\to\{0,1,*\}$, with an approximation ratio better than $c$.
\end{ntcorollary}
Let us now prove that, assuming $\operatorname P\ne\operatorname{NP}$, there is no polynomial-time computable approximate characterization of the communication complexity of a partial Boolean function. This is captured by a more general result (Theorem35, stated at the start of this item), which is an immediate consequence of the reduction presented in Section5.2 and Theorem21.
\begin{proof}[Author proof of Theorem35]
Suppose that functions $\lambda,\eta$, and $r$ exist as described in the statement of the theorem. Let $L$ be any language in $\operatorname{NP}$, and assume that $k$ is a sufficiently large constant. For any $x\in\{0,1\}^n$, let $G_x$ be the corresponding graph given by Theorem21. Then let $f_x=f_{G_x}$ be as defined in Section5.2 above, and let $M_x$ be the communication matrix of $f_x$. If $x\in L$, we then have by Theorem32 that $k\geq\chi(G_x)\geq\frac12L(f_x)$, and because the protocol for $f_x$ that witnesses this fact is balanced, we get $1+\log k\geq D(f_x)$. Hence $\eta(1+\log k)\geq\eta(D(f_x))\geq r(M_x)$. If, on the other hand,$y\notin L$, then we have $2^{0.99k}<\binom k{\lfloor k/2\rfloor}<\chi(G_y)=C_1(f_y)\leq L(f_y)$, and thus $0.99k<D(f_y)$. Then $\lambda(0.99k)\leq\lambda(D(f_y))\leq r(M_y)$. But since $\lambda(0.99k)>\eta(1+\log k)$ using our assumptions on these functions and on $k$, it follows that $r(M_y)>r(M_x)$. As a consequence, the polynomial time function $r$ can be used to distinguish positive and negative instances of $L$. Since $L$ is an arbitrary language in $\operatorname{NP}$, we get that $\operatorname P=\operatorname{NP}$. This completes the proof.
\end{proof}

\begin{thebibliography}{82}
\bibitem{ref1}Misha Alekhnovich, Mark Braverman, Vitaly Feldman, Adam R. Klivans, and Toniann Pitassi. The complexity of properly learning simple concept classes. Journal of Computer and System Sciences, 74(1):16–34, 2008.
\bibitem{ref2}Eric Allender and Bireswar Das. Zero knowledge and circuit minimization. In International Symposium on Mathematical Foundations of Computer Science (MFCS), pages 25–32, 2014.
\bibitem{ref3}Eric Allender, Joshua A. Grochow, Dieter van Melkebeek, Cristopher Moore, and Andrew Morgan. Minimum circuit size, graph isomorphism, and related problems. SIAM Journal on Computing, 47(4):1339–1372, 2018.
\bibitem{ref4}Eric Allender, Lisa Hellerstein, Paul McCabe, Toniann Pitassi, and Michael E. Saks. Minimizing disjunctive normal form formulas and AC0 circuits given a truth table. SIAM Journal on Computing, 38(1):63–84, 2008.
\bibitem{ref5}Eric Allender and Shuichi Hirahara. New insights on the (non-)hardness of circuit minimization and related problems. In International Symposium on Mathematical Foundations of Computer Science (MFCS), pages 54:1–54:14, 2017.
\bibitem{ref6}Eric Allender, Dhiraj Holden, and Valentine Kabanets. The minimum oracle circuit size problem. Computational Complexity, 26(2):469–496, 2017.
\bibitem{ref7}Eric Allender, Rahul Ilango, and Neekon Vafa. The non-hardness of approximating circuit size. In International Computer Science Symposium in Russia (CSR), pages 13–24, 2019.
\bibitem{ref8}Anurag Anshu, Naresh Goud Boddu, and Dave Touchette. Quantum log-approximate-rank conjecture is also false. arXiv, 2018. arXiv:1811.10525.
\bibitem{ref9}Benny Applebaum, Boaz Barak, and David Xiao. On basing lower-bounds for learning on worst-case assumptions. In Symposium on Foundations of Computer Science (FOCS), pages 211–220, 2008.
\bibitem{ref10}Theodore P. Baker, John Gill, and Robert Solovay. Relativizations of the P =? NP question. SIAM Journal on Computing, 4(4):431–442, 1975.
\bibitem{ref11}Avrim Blum, Merrick L. Furst, Michael J., and Richard J. Lipton. Cryptographic primitives based on hard learning problems. In International Cryptology Conference (CRYPTO), pages 278–291, 1993.
\bibitem{ref12}Avrim L. Blum and Ronald L. Rivest. Training a 3-node neural network is np-complete. Neural Networks, 5(1):117–127, 1992.
\bibitem{ref13}Anselm Blumer, Andrzej Ehrenfeucht, David Haussler, and Manfred K. Warmuth. Occam’s razor. Information Processing Letters, 24(6):377–380, 1987.
\bibitem{ref14}Andrej Bogdanov and Luca Trevisan. Average-case complexity. Foundations and Trends in Theoretical Computer Science, 2(1), 2006.
\bibitem{ref15}Dan Boneh and Richard J. Lipton. Amplification of weak learning under the uniform distribu- tion. In Conference on Learning Theory (COLT), pages 347–351, 1993.
\bibitem{ref16}Joan Boyar, Philip Matthews, and René Peralta. On the shortest linear straight-line program for computing linear forms. In International Symposium on Mathematical Foundations of Computer Science (MFCS), pages 168–179, 2008.
\bibitem{ref17}Marco L. Carmosino, Russell Impagliazzo, Valentine Kabanets, and Antonina Kolokolova. Learning algorithms from natural proofs. In Conference on Computational Complexity (CCC), pages 10:1–10:24, 2016.
\bibitem{ref18}Amit Chakrabarti and Oded Regev. An optimal lower bound on the communication complexity of gap-Hamming-distance. SIAM Journal on Computing, 41(5):1299–1317, 2012. doi:10. 1137/120861072.
\bibitem{ref19}Arkadev Chattopadhyay, Nikhil S Mande, and Suhail Sherif. The log-approximate-rank conjecture is false. In Symposium on Theory of Computing (STOC), pages 42–53. ACM, 2019.
\bibitem{ref20}Mahdi Cheraghchi, Valentine Kabanets, Zhenjian Lu, and Dimitrios Myrisiotis. Circuit lower bounds for MCSP from local pseudorandom generators. Electronic Colloquium on Computational Complexity (ECCC), 26:22, 2019.
\bibitem{ref21}Gil Cohen and Igor Shinkar. The complexity of DNF of parities. In Innovations in Theoretical Computer Science (ITCS), pages 47–58, 2016.
\bibitem{ref22}Sebastian Czort. The complexity of minimizing disjunctive normal form formulas. Master’s Thesis, University of Aarhus, 1999.
\bibitem{ref23}Uriel Feige. A threshold of ln n for approximating set cover. Journal of the ACM, 45(4):634–652, 1998.
\bibitem{ref24}Uriel Feige and Joe Kilian. Zero knowledge and the chromatic number. Journal of Computer and System Sciences, 57(2):187–199, 1998.
\bibitem{ref25}Magnus Gausdal Find, Alexander Golovnev, Edward A. Hirsch, and Alexander S. Kulikov. A better-than-3n lower bound for the circuit complexity of an explicit function. In Symposium on Foundations of Computer Science (FOCS), pages 89–98, 2016.
\bibitem{ref26}Michael R. Garey and David S. Johnson. Computers and Intractability: A Guide to the Theory of NP-Completeness. W. H. Freeman, 1979.
\bibitem{ref27}Alexander Golovnev, Rahul Ilango, Russell Impagliazzo, Valentine Kabanets, Antonina Koloko- lova, and Avishay Tal. AC0 [p] lower bounds against MCSP via the coin problem. Electronic Colloquium on Computational Complexity (ECCC), 26:18, 2019.
\bibitem{ref28}Mika Göös and Toniann Pitassi. Communication lower bounds via critical block sensitivity. SIAM Journal on Computing, 47(5):1778–1806, 2018. doi:10.1137/16M1082007.
\bibitem{ref29}Gary D. Hachtel and Fabio Somenzi. Logic synthesis and verification algorithms. Springer, 2006.
\bibitem{ref30}Thomas R. Hancock, Tao Jiang, Ming Li, and John Tromp. Lower bounds on learning decision lists and trees. Information and Computation, 126(2):114–122, 1996.
\bibitem{ref31}Johan Hastad. Clique is hard to approximate within $n^{1-\varepsilon}$ . In Symposium on Foundations of Computer Science (FOCS), pages 627–636, 1996.
\bibitem{ref32}Shuichi Hirahara. Non-black-box worst-case to average-case reductions within NP. In Sym- posium on Foundations of Computer Science (FOCS), pages 247–258, 2018.
\bibitem{ref33}Shuichi Hirahara, Igor Carboni Oliveira, and Rahul Santhanam. NP-hardness of minimum circuit size problem for OR-AND-MOD circuits. In Computational Complexity Conference (CCC), pages 5:1–5:31, 2018.
\bibitem{ref34}Shuichi Hirahara and Rahul Santhanam. On the average-case complexity of MCSP and its variants. In Computational Complexity Conference (CCC), pages 7:1–7:20, 2017.
\bibitem{ref35}Shuichi Hirahara and Osamu Watanabe. Limits of minimum circuit size problem as oracle. In Conference on Computational Complexity (CCC), pages 18:1–18:20, 2016.
\bibitem{ref36}John M. Hitchcock and Aduri Pavan. On the NP-completeness of the minimum circuit size problem. In Conference on Foundation of Software Technology and Theoretical Computer Science (FSTTCS), pages 236–245, 2015.
\bibitem{ref37}Sangxia Huang. Improved hardness of approximating chromatic number. In International Con- ference on Approximation Algorithms for Combinatorial Optimization Problems (APPROX), pages 233–243, 2013.
\bibitem{ref38}Rahul Ilango. AC0 [p] lower bounds and NP-hardness for variants of MCSP. Electronic Colloquium on Computational Complexity (ECCC), 26:21, 2019.
\bibitem{ref39}Russell Impagliazzo, Valentine Kabanets, and Ilya Volkovich. The Power of Natural Properties as Oracles. In Computational Complexity Conference (CCC), volume 102, pages 7:1–7:20, 2018.
\bibitem{ref40}Kazuo Iwama and Hiroki Morizumi. An explicit lower bound of 5n − o(n) for boolean circuits. In International Symposium on Mathematical Foundations of Computer Science (MFCS), pages 353–364, 2002.
\bibitem{ref41}J. Stephen Judd. Learning in networks is hard. In International Conference on Neural Networks (ICNN), volume 2, pages 685–692, 1987.
\bibitem{ref42}Stasys Jukna. Boolean function complexity: advances and frontiers, volume 27. Springer, 2012.
\bibitem{ref43}Valentine Kabanets and Jin-yi Cai. Circuit minimization problem. In Symposium on Theory of Computing (STOC), pages 73–79, 2000.
\bibitem{ref44}Maurice Karnaugh. The map method for synthesis of combinational logic circuits. Transactions of the American Institute of Electrical Engineers, Part I: Communication and Electronics, 72(5):593–599, 1953.
\bibitem{ref45}Michael J. Kearns and Leslie G. Valiant. Cryptographic limitations on learning boolean formulae and finite automata. J. ACM, 41(1):67–95, 1994. doi:10.1145/174644.174647.
\bibitem{ref46}Michael J. Kearns and Umesh V. Vazirani. An Introduction to Computational Learning Theory. MIT Press, 1994.
\bibitem{ref47}Subhash Khot. Improved inapproximability results for maxclique, chromatic number and approximate graph coloring. In Symposium on Foundations of Computer Science (FOCS), pages 600–609, 2001.
\bibitem{ref48}Subhash Khot and Rishi Saket. Hardness of minimizing and learning DNF expressions. In Symposium on Foundations of Computer Science (FOCS), pages 231–240, 2008.
\bibitem{ref49}Ker-I Ko. On the complexity of learning minimum time-bounded Turing machines. SIAM Journal on Computing, 20(5):962–986, 1990.
\bibitem{ref50}Pravesh K. Kothari and Roi Livni. Improper learning by refuting. In Innovations in Theoretical Computer Science (ITCS), pages 55:1–55:10, 2018.
\bibitem{ref51}Jan Krajícek. Forcing with Random Variables and Proof Complexity. Cambridge University Press, 2011.
\bibitem{ref52}Eyal Kushilevitz. Communication complexity. In Advances in Computers, volume 44, pages 331–360. Elsevier, 1997.
\bibitem{ref53}Eyal Kushilevitz and Noam Nisan. Communication Complexity. Cambridge University Press, 1997.
\bibitem{ref54}Eyal Kushilevitz and Enav Weinreb. On the complexity of communication complexity. In Symposium on Theory of Computing (STOC), pages 465–474, 2009.
\bibitem{ref55}Leonid Levin. Universal sequential search problems. Problemy Peredachi Informatsii, 9(3):115– 116, 1973.
\bibitem{ref56}László Lovász and Michael Saks. Lattices, mobius functions and communications complexity. In Symposium on Foundations of Computer Science (FOCS), pages 81–90, 1988.
\bibitem{ref57}Carsten Lund and Mihalis Yannakakis. On the hardness of approximating minimization problems. Journal of the ACM, 41(5):960–981, 1994.
\bibitem{ref58}William J. Masek. Some NP-complete set covering problems. Unpublished Manuscript, 1979.
\bibitem{ref59}Edward J. McCluskey. Introduction to the theory of switching circuits. McGraw-Hill, 1965.
\bibitem{ref60}Dylan M. McKay, Cody D. Murray, and R. Ryan Williams. Weak lower bounds on resource- bounded compression imply strong separations of complexity classes. In Symposium on Theory of Computing (STOC), 2019.
\bibitem{ref61}Moritz Müller and Ján Pich. Feasibly constructive proofs of succinct weak circuit lower bounds. Electronic Colloquium on Computational Complexity (ECCC), 24:144, 2017.
\bibitem{ref62}Cody D. Murray and Richard Ryan Williams. On the (non) NP-hardness of computing circuit complexity. In Conference on Computational Complexity (CCC), pages 365–380, 2015.
\bibitem{ref63}Igor Carboni Oliveira, Ján Pich, and Rahul Santhanam. Hardness magnification near state-of- the-art lower bounds. In Conference on Computational Complexity (CCC), 2019.
\bibitem{ref64}Igor Carboni Oliveira and Rahul Santhanam. Conspiracies between learning algorithms, circuit lower bounds, and pseudorandomness. In Computational Complexity Conference (CCC), pages 18:1–18:49, 2017.
\bibitem{ref65}Igor Carboni Oliveira and Rahul Santhanam. Hardness magnification for natural problems. In Symposium on Foundations of Computer Science (FOCS), pages 65–76, 2018.
\bibitem{ref66}James Orlin. Contentment in graph theory: covering graphs with cliques. Indagationes Mathematicae, 80(5):406–424, 1977.
\bibitem{ref67}Leonard Pitt and Leslie G. Valiant. Computational limitations on learning from examples. Journal of the ACM, 35(4):965–984, 1988.
\bibitem{ref68}Pavel Pudlák, Vojtech Rödl, and Petr Savický. Graph complexity. Acta Informatica, 25(5):515– 535, 1988.
\bibitem{ref69}Netanel Raviv.     Truth table minimization of computational models.           CoRR/arXiv, abs/1306.3766, 2013. arXiv:1306.3766.
\bibitem{ref70}Alexander A. Razborov and Steven Rudich. Natural proofs. Journal of Computer and System Sciences, 55(1):24–35, 1997.
\bibitem{ref71}Ronald L. Rivest. Learning decision lists. Machine Learning, 2(3):229–246, 1987.
\bibitem{ref72}J. Paul Roth and Richard M. Karp. Minimization over boolean graphs. IBM Journal of Research and Development, 6(2):227–238, 1962.
\bibitem{ref73}Christoph Scholl. Functional decomposition with applications to FPGA synthesis. Kluwer Academic Publishers, 2001.
\bibitem{ref74}Makrand Sinha and Ronald de Wolf. Exponential separation between quantum communication and logarithm of approximate rank. arXiv, 2018. arXiv:1811.10090.
\bibitem{ref75}Boris A Trakhtenbrot. A survey of Russian approaches to perebor (brute-force search) algorithms. Annals of the History of Computing, 6(4):384–400, 1984.
\bibitem{ref76}Luca Trevisan. Non-approximability results for optimization problems on bounded degree instances. In Symposium on Theory of Computing (STOC), pages 453–461, 2001.
\bibitem{ref77}Christopher Umans, Tiziano Villa, and Alberto L. Sangiovanni-Vincentelli. Complexity of two-level logic minimization. IEEE Trans. on CAD of Integrated Circuits and Systems, 25(7):1230–1246, 2006.
\bibitem{ref78}Salil P. Vadhan. On learning vs. refutation. In Conference on Learning Theory (COLT), pages 1835–1848, 2017.
\bibitem{ref79}Marcin Wrochna and Stanislav Živnỳ. Improved hardness for H-colourings of G-colourable graphs. arXiv, 2019. arXiv:1907.00872.
\bibitem{ref80}Andrew Chi-Chih Yao. Some complexity questions related to distributive computing (pre- liminary report). In Symposium on Theory of Computing (STOC), pages 209–213, 1979. doi:10.1145/800135.804414.
\bibitem{ref81}Andrew Chi-Chih Yao. Theory and applications of trapdoor functions (extended abstract). In Symposium on Foundations of Computer Science (FOCS), pages 80–91, 1982.
\bibitem{ref82}David Zuckerman. Linear degree extractors and the inapproximability of max clique and chromatic number. In Symposium on Theory of Computing (STOC), pages 681–690, 2006.
\end{thebibliography}
\end{document}
